% Данный файл распространяетсяя под лицензией CC BY 4.0. Текст лицензии размещён на https://creativecommons.org/licenses/by/4.0/
% (c) Кафедра системного программирования, 2025

\documentclass[a4paper]{article}

\usepackage[a4paper, top=8mm, bottom=8mm, left=8mm, right=8mm]{geometry}

\usepackage{polyglossia}
\setdefaultlanguage[babelshorthands=true]{russian}
\setotherlanguage{english}

\usepackage{fontspec}
\setmainfont{FreeSerif}
\newfontfamily{\russianfonttt}[Scale=0.7]{DejaVuSansMono}

\usepackage[tiny, compact]{titlesec}

\usepackage{titling}
\setlength{\droptitle}{-1cm}
\pretitle{\begin{center}\begin{bfseries}\Large}
\posttitle{\par\end{bfseries}\end{center}}
\preauthor{\begin{center}\normalsize}
\postauthor{\par\end{center}\vspace{-1.8cm}}

\usepackage{hyperref}
\usepackage{bookmark}
\usepackage{csquotes}
% Сюда можно дописывать нужные вам пакеты.

\title{Улучшение методов дедупликации SBC}

\author{Никитин Артём Денисович}

% Дата много места занимает, её лучше не указывать.
\date{}

\begin{document}

\maketitle

\begin{flushright}
    Группа: \emph{24.Б44-мм}

    Кафедра: \emph Системное программирование

    % Степень/звание и должность научника мы лучше вас знаем, их можно не указывать.
    Научный руководитель: \emph{Гориховский Вячеслав Игоревич}

    % А вот про консультанта укажите официальное место работы, с "ООО" или чем-то таким, и должностью желательно по трудовой, не просто "techlead". Выясните у консультанта.
    Консультант: \emph{Дмитриевцев Алексей Сергеевич}
    
    Номер семестра практики: \emph{3}
\end{flushright} 


\section{Постановка задачи}

% Буквально парой предложений ввести в контекст, обязательно сформулировать цель работы. Если работа делается где-то глубоко внутри проекта, на введение можно потратить больше места, чтобы подвести читателя от сути проекта в целом к тому, что конкретно вам надо было сделать. Только без воды, никакого "в современном мире"!
Работа в интересах компании YADRO в рамках проекта ChunkFS. \textbf{Целью работы является улучшить методы дедупликации SBC, а также создать механизмы их анализа в файловой системе ChunkFS}

\textbf{Дедупликация данных} — это алгоритм сжатия, нацеленный на выявление
идентичных блоков данных и их удаление.

\textbf{Similarity Based Chunking (SBC)} — метод дедупликации, при
котором данные сначала разбиваются на блоки (чанки), а затем с помощью
хеширования определяется степень их схожести. Похожие чанки объединяются
в кластеры. Сжатие обеспечивается алгоритмами дельта-кодирования, такими
как Zdelta, Xdelta, Gdelta, Ddelta и Edelta. Эти кодировщики итерируются по
кластеру, оставляя в нём только один родительский чанк и инструкции для
восстановления остальных похожих чанков.

% Это пояснение актуальности. Тут обязательно надо обрисовать полезность вашей работы для конечного пользователя (и заодно как-то описать, кто такой ваш конечный пользователь --- может, это соседняя команда). Даже если вы делаете работу "глубоко внутри", выясните у консультанта, чем это людям поможет. И опишите, кто выступал в роли "заказчика" вашей работы, т.е. кто и зачем предложил задачу. Можно без конкретных имён, но названия организаций вполне желательны.
Результаты работы необходимы для качественного бенчмаркинга SBC-алгоритма. Исследования проводятся в целях увеличения коэффициента дедупликации.

\section{Описание предлагаемого решения}

% Раздел пишется в свободной форме, с минимумом технических подробностей (можно приводить ссылки на более подробные документы).

Xdelta\footnote{Dimitre Trendafilov, Nasir Memon, and Torsten Suel. Zdelta: An efficient delta compression tool. Technical report,
Department of Computer and Information Science at Polytechnic University, 2002.} — алгоритм дельта-кодирования использующий скользящее окно для вычисления хеша Adler32.

Gdelta\footnote{H. Tan, Z. Zhang, X. Zou, Q. Liao, and W. Xia, “Exploring the potential of fast delta encoding: Marching to a higher compression ratio” in Proc. IEEE Conf., 2020.} — использует для хеширования не Adler32, а gear hash. (Было реализовано ранее)

Ddelta\footnote{Wen Xia, Hong Jiang, Dan Feng, Lei Tian, Min Fu, and Yukun Zhou. Ddelta: A deduplication-inspired fast delta compression approach. Performance Evaluation,
79:258–272, 2014.} — алгоритм использующий content define chunking (CDC) для разбиения чанка и вычисления хеша.

Edelta\footnote{Wen Xia, Chunguang Li, Hong Jiang, Dan Feng, Yu Hua, Leihua Qin, and Yucheng Zhang. Edelta: A word-enlarging based fast delta compression approach. In the
7th USENIX conference on Hot Topics in Storage and File Systems (HotStorage’15), pages 1–5, Santa Clara, CA, 2015. USENIX Association.} — оптимизация ddelta, использующая эвристику: если данные в целевом чанке и родительском совпали, то данные находящиеся поблизости тоже совпадут.

Zdelta\footnote{Haoliang Tan , Zhiyuan Zhang , Xiangyu Zou , Qing Liao, Wen Xia. Exploring the Potential of Fast Delta Encoding: Marching to a Higher Compression Ratio. — 2021.} — алгоритм использующий три указателя для умного вычисления смещения copy инструкций.

Алгоритмы интегрированы в ChunkFS — файловую систему для бенчмаркинга алгоритмов
дедупликации.

Помимо этих алгоритмов мной была реализована функциональность позволяющая собирать данные о этапе кластеризации SBC-алгоритма.

Ранее существующие решения не интегрированы в файловые системы бенчмаркинга методов дедупликации с открытым и исходным кодом.

% Вот это важно, опишите кратко, что использовали
Реализация решения осуществлялась с использованием языка Rust и файловой системы ChunkFS.

\section{Эксперименты}

% Название секции можно поменять на более подходящее вашей работе

% Название этой секции --- либо \enquote{Апробация}, либо \enquote{эксперименты}  (если содержательно есть и то и то, можно две отдельные секции сделать). Тут опишите, как проводилась проверка того, что получилось что-то адекватное. Если есть отзывы пользователей, отлично (особенно если они в структурированном виде, типа анкеты SUS). Если есть отзывы команды, консультанта и т.п., хорошо. Если есть только тесты. то хотя бы их тут опишите.

% Если работа подразумевала замеры чего-то (например, производительности), приведите тут итоговые результаты и выводы, и дайте ссылку на таблицы с данными. Не забывайте про матстат (приводите матожидание и среднеквадратичное отклонение, не просто результат замера).
\textbf{Исследовательский вопрос:} как коэффициент дедупликации для каждого кодировщика зависит от CDC?

\textbf{Метрики:} \url{https://docs.google.com/spreadsheets/d/152O0nph2a-DgmPHKsB0gL2AIYR7s2dLHWC19rbjdrLo/edit?hl=ru&gid=0#gid=0}

\textbf{Вывод:} Наилучший коэффициент дедупликации для всех кодировщиков был получен на Ultra CDC (прирост около 10 процентов). Ранее при анализе других кодировщиков были получены схожие результаты\footnote{Щербаков М.А. Гориховский В.И. Повышение качества 
дедупликации с помощью Similarty Based Chunking 2025. –– \url{https://www.elibrary.ru/item.asp?id=81439119&pff=1}}. Можно сделать предположение что Ultra это самый подходящий CDC для работы с SBC. Для всех CDC Xdelta даёт наибольший коэффициент дедупликации.

\section{Заключение}

В ходе выполнения работы получены следующие \textbf{результаты}:

% ниже списком перечисляете то, что выносится как результат практики.

\begin{itemize}
    \item Реализованы алгоритмы xdelta, ddelta, zdelta и
edelta для SBC scrubber
    \item Добавлена возможность собирать статистику
на этапе кластеризации SBC алгоритма
    \item Установлена зависимость коэффициента дедупликации от CDC для различных дельта-кодировщиков
\end{itemize}

Ссылки на пуллреквесты: 
\begin{itemize}
    \item Clustering measurements in ChunkFS: \url{https://github.com/Piletskii-Oleg/chunkfs/pull/9}
    \item Clustering measurements in SBC scrubber: \url{https://github.com/maxscherbakov/sbc_algorithm/pull/13}
    \item Xdelta: \url{https://github.com/maxscherbakov/sbc_algorithm/pull/9}
    \item Ddelta/Edelta: \url{https://github.com/maxscherbakov/sbc_algorithm/pull/12}
    \item Zdelta: \url{https://github.com/maxscherbakov/sbc_algorithm/pull/10}
\end{itemize}

Техническая документация: \url{https://github.com/ArtemNikit1n/sbc_algorithm/blob/main/src/encoder/EncoderDocument.md}

В будущем планируется исследовать взаимосвязь коэффициента дедупликации и формы кластера на различных датасетах.

\end{document}
